\section{Introduction}
\vspace{-0.3cm}

Large Language Model (LLM) inference is a rapidly growing workload. It has two phases~\cite{flashdecoding++}: (i) the \textit{prefill phase}, which processes input tokens (the prompt) and spends most of its cycles on General Matrix Multiply (GEMM); and (ii) the \textit{decode phase}, which generates tokens one by one in an autoregressive manner, primarily performing General Matrix-Vector Product (GEMV). Decode requires repeatedly loading the entire LLM model into on-chip memory, with GEMV dominating its cycles. Since LLMs generate many tokens, especially in the test-time scaling scenario, such as the OpenAI-o1/o3~\cite{openaio1,openaio3} and DeepSeek-R1~\cite{deepseekr1}, inference is constrained by GEMV latency, making it inherently memory-bandwidth-bound.

To address memory bandwidth bottlenecks, AI accelerators are increasingly adopting system-on-wafer integration~\cite{tsmc-advanced-packaging}. This approach scales chip area to a full wafer, up to 100$\times$ larger than a typical GPU die, enabling significantly more on-chip cores, memory, and bandwidth. Examples include Cerebras WSE~\cite{wse} and upcoming Tesla Dojo~\cite{dojo}. The Cerebras WSE-2, for instance, integrates 850,000 cores with 40GB of on-chip memory, 1,000$\times$ more than GPUs, and provides 22PB/s memory bandwidth, 7,000$\times$ higher than GPUs. TSMC anticipates widespread adoption of system-on-wafer integration, citing advantages in performance, energy-efficient die-to-die communication, and cost reduction. IEEE similarly forecasts a surge in wafer-scale computing by 2027~\cite{tsmc-advanced-packaging}. Wafer-scale accelerators are already seeing real-world deployment, particularly in model serving. In February 2025, Mixtral and Perplexity adopted wafer-scale chips, achieving cost parity with GPUs in terms of tokens per dollar~\cite{perplexity,mistra}. G42 now operates data centers fully outfitted with wafer-scale accelerators, and Cerebras has secured major commercial contracts~\cite{g42}.

Unlocking the potential of wafer-scale accelerators is challenging because current LLM systems rely on \emph{shared memory architectures} typical of GPUs and TPUs. Wafer-scale accelerators, however, adopt \emph{network-on-chip} (NoC) designs that interconnect millions of cores with local memory in a \emph{massive-scale, mesh-based memory architecture}. This architecture far exceeds the scale of on-chip crossbars (e.g., one-hop NUMA such as GraphCore IPU), multi-socket NUMA~\cite{amd2023numa}, and high-density AI clusters (hundreds of GPUs per pod)~\cite{tpu}. Without fully addressing this fundamental shift in memory architecture, directly applying designs from state-of-the-art systems like T10~\cite{t10} and Ladder~\cite{ladder} to wafer-scale devices often results in extremely poor performance.

To address these challenges, we propose a \emph{device model} that captures the critical hardware properties of wafer-scale accelerators, highlighting key differences from shared-memory devices. This model enables us to evaluate current LLM inference design principles, identify non-compliant areas, and pinpoint where new approaches are required. Guided by this model, we can achieve an ambitious system design: \emph{running complete LLM inference on a single chip}, minimizing costly off-chip communication and maximizing on-chip memory bandwidth utilization.

The above idea motivates \sys, the first wafer-scale LLM inference system, yielding several contributions:

\mypar{(1) Device model for wafer-scale accelerators}
We propose the PLMR model\footnote{PLMR model can be pronounced as “Plummer” }, which captures the key hardware properties of wafer-scale accelerators:
(i) Massive \textbf{P}arallel cores (P): Millions of cores can be integrated on a large wafer, requiring systems to effectively partition LLMs and their operations.
(ii) Highly non-uniform memory access \textbf{L}atency (L): Inter-core data access exhibits significant variation, with latency differences up to 1,000$\times$, necessitating the system to mitigate this.
(iii) Constrained per-core local \textbf{M}emory (M): Each core has limited memory (tens of KBs to several MBs), requiring efficient memory usage.
(iv) Limited hardware-assisted \textbf{R}outing (R): The NoC routing hardware is constrained by the area size per core and can only support a limited number of routing paths, e.g., less than $2^5$ on Cerebras WSE-2.

\mypar{(2) Wafer-scale LLM parallelism}
We propose an effective, PLMR-compliant LLM parallelism policy for wafer-scale accelerators. %, fully compliant with the PLMR model. 
In the prefill phase, we design fine-grained partitioning to achieve million-core parallelism. For the decode phase, where tensor dimensions are insufficient for partitioning, we design fine-grained replication to enable parallelism with minimal communication costs. As a result, \sys achieves larger-scale and finer-grained parallelism (satisfying P in PLMR) than GPU-based approaches. Additionally, we replace conventional GPU-based GEMM and GEMV operators %in existing LLM models 
with new algorithm designed for the PLMR model (satisfying L, M, and R) and propose tensor placement strategies that eliminate matrix transpositions, which are costly with a mesh NoC (satisfying L). 

We also designed a scalable KV-cache management method for wafer-scale devices. This approach features a novel KV cache shift method to ensure balanced core usage (satisfying P and M), avoiding skewed utilization of cores caused by KV cache concatenation methods common on GPUs.

\mypar{(3) Wafer-scale GEMM}
We propose \gemm, a scalable GEMM algorithm for wafer-scale devices, accelerating the prefill phase.
Unlike conventional distributed GEMM algorithms, \gemm achieves full PLMR compliance by leveraging two key operations: \emph{cyclic shifting} and \emph{interleaving}. Cyclic shifting ensures algorithm correctness while maintaining bounded usage of local memory (satisfying M). The interleaving operation minimizes communication latency in the mesh NoC, effectively reducing the overhead of highly non-uniform memory latency (satisfying L) and routing resources (satisfying R).

\mypar{(4) Wafer-scale GEMV} 
We propose \gemv, a scalable GEMV algorithm for wafer-scale devices, accelerating the decode phase. Unlike existing GEMV implementations, \gemv uses a novel \emph{K-tree allreduce} algorithm to aggregate local GEMV results across massive cores. This algorithm ensures routing resource usage meets the hardware's limitation (satisfying R) and reduces communication latency (satisfying L).

\vspace{0.2cm}
\noindent
We implemented \sys on the Cerebras WSE engine using approximately 7,000 lines of CSL (a C-like programming language) for LLM parallelism, \gemm, and \gemv, and 2,000 lines of Python for loading LLM checkpoints, launching inference, and execution parallelism policy.

We conducted end-to-end LLM inference experiments with various models, including full LLaMA3-8B and LLaMA2-13B, as well as subsets of layers of CodeLLaMA-34B, and QWen2-72B. By combining wafer-scale LLM parallelism, GEMM and GEMV, \sys outperforms state-of-the-art (SOTA) systems:
(i) 100-200$\times$ faster than T10~\cite{t10}, the SOTA system for massive cores with a distributed on-chip memory architecture, and
(ii) 200-400$\times$ faster than Ladder~\cite{ladder}, the SOTA system for shared-memory architectures.

Micro-benchmarks further show that \gemm is 2-3$\times$ faster than SUMMA~\cite{summa}, the default optimized GEMM for Cerebras WSE, and Cannon~\cite{cannon}, the SOTA GEMM for supercomputers with large-scale mesh architectures. \gemv achieves 4-8$\times$ speedups over Cerebras’s optimized GEMV, and 606$\times$ than a single A100 GPU. Additionally, \sys’s cache shift method is up to 400$\times$ more scalable than the KV cache SOTA on GPUs, such as PagedAttention~\cite{vllm}.

Combined, \sys (on Cerebras WSE-2) outperforms SGLang (on single A100) by 30-40$\times$. Compared to the optimal performance of SGLang on A100 multi-GPUs connected with NVLink and RDMA, \sys delivers a 10-20$\times$ faster end-to-end speed and is 2.5$\times$ more energy-efficient. The reduced gains from GEMV to LLM are due to current limitations in software, hardware, and existing LLM model designs. We anticipate stronger performance as wafer-scale AI computing matures and these limitations are addressed.